\documentclass[10pt,a4paper]{amsart}
\usepackage[margin=19mm]{geometry}
\usepackage{amsmath,amssymb,mathtools}
\usepackage[scr=rsfs,cal=euler]{mathalfa}
\usepackage{xfrac}
\usepackage[unicode,hidelinks,bookmarks=false]{hyperref}
\newcommand{\ie}{i.e.\ }
\newcommand{\cf}{cf.\ }
\newcommand{\resp}{respectively\ }
\newcommand{\Subsection}{\subsection}
\providecommand{\nobreakdash}{\nobreak\hbox{-}}
\setlength{\emergencystretch}{2em}
\multlinegap0pt
\usepackage{bm}
\usepackage{bbm}
\usepackage{dsfont}
\usepackage{xspace}
\usepackage{cancel}
\usepackage{mathtools}
\usepackage{amsthm}
\makeatletter
\@ifundefined{theorem}{\newtheorem{theorem}{Theorem}}{}
\@ifundefined{proposition}{\newtheorem{proposition}[theorem]{Proposition}}{}
\makeatother
\renewcommand{\subset}{\subseteq}
\title[Relative bi-exactness and structural results for graph-wreat]{Relative bi-exactness and structural results for graph-wreath product von Neumann algebras}
\author{Taisuke Hoshino}
\date{Source: arXiv:2601.18185v1 (CC BY 4.0)}
\begin{document}
\maketitle

\noindent\emph{Source and licence.} Relative bi-exactness and structural results for graph-wreath product von Neumann algebras. Version arXiv:2601.18185v1 (CC BY 4.0), distributed under the Creative Commons Attribution 4.0 International licence, as recorded in the arXiv licence metadata for this exact version and shown on the abstract page https://arxiv.org/abs/2601.18185v1. Full rights notice: licenses/arxiv-2601.18185-CC-BY-4.0.txt.\par\smallskip

\noindent\textit{Editorial scope.} Counted unit: the Proposition whose source label is \texttt{None} in the article (source0.tex lines 829--848, source label \texttt{None}), delivered together with the article's own complete original proof of it (source0.tex lines 850--925, one \texttt{proof} environment of 76 lines). No mathematics is written, repaired or re-derived: statement and proof are the article's own text line for line. Required context retained uncounted: none. Every definition, display and label of those lines is retained unchanged. Omitted from this package and not redistributed: 1--828, 849--849, 926--1607. Nothing inside a retained excerpt is omitted or altered: every retained body is delivered exactly as the article's lines are written, so no omission receipt of any kind accompanies this package. Citations: the retained excerpts carry no citation command at all, so no bibliography entry of the article is transcribed and no reference is redistributed. Reference adaptation: none was needed and none was made; every reference inside the delivered unit resolves inside it, so no unresolved reference or citation is produced. Printed slips kept literal: no printed word, symbol, hypothesis or number of the counted statement or of its proof was altered. Layout and class: the article's own class is \texttt{amsart} with packages including amssymb, amsmath, amsfonts, amsthm, mathtools, mathrsfs, xspace, tikz-cd, enumitem, graphicx; the wrapper is the lane's pinned \texttt{amsart} editorial wrapper at 10pt on A4, which restates the theorem-like environments the retained text needs and adds the article's own macro definitions that the retained text needs (\texttt{\textbackslash subset}), with extra wrapper packages (bm, bbm, dsfont, xspace, cancel, mathtools). The delivered numbering is the unit's own. Assets: no retained span contains a figure environment, a graphics-inclusion command, a PDF-page inclusion, a table or a TikZ picture; no image file is copied into this package, the package declares no asset dependency and no image file is redistributed with it. Licence: version arXiv:2601.18185v1 of this article is distributed under the Creative Commons Attribution 4.0 International licence, as recorded in the arXiv licence metadata for this exact version and shown on the abstract page https://arxiv.org/abs/2601.18185v1. A full rights notice is delivered at licenses/arxiv-2601.18185-CC-BY-4.0.txt. Characters: every retained excerpt is pure ASCII, so the delivered engine, XeLaTeX with its default Latin Modern text font, reports no missing character.
\begin{proposition}
The following holds for $z=(h=h_1h_2\cdots h_m, g)\in H_{\Gamma}\rtimes G$ and
$\mathcal{G}=\{ A(E, F, n) \mid E\subset V\Gamma, F\subset H, n\} \cup\{ H_{\mathrm{Lk}v}\mid v\in V\Gamma\}$; 
\begin{enumerate}
\item $|z|_f\to \infty$ when $z\to \infty/\mathcal{G}$. 

\item $\dfrac{|\mathop{supp}h|}{|z|_f} \to 0$ when $z\to \infty/\mathcal{G}$. 

\item For any $v\in V\Gamma$ and $x\in H_v$, 
\begin{align}
\| m(xz)-m(z) \|_1, \| m(zx)-m(z) \|_1 \leqq |x|_f.  
\end{align}
In particular, $\dfrac{\| m(xz)-m(z) \|_1}{|z|_f}, \dfrac{\| m(zx)-m(z) \|_1}{|z|_f} \to 0$ when $z\to \infty/\mathcal{G}$.
\item For any $k\in G$, 
\begin{align}
\| m(kz)-k.m(z) \|_1, \| m(zk)-m(z) \|_1\leqq |k|_G|\mathop{supp}h|. 
\end{align}
In particular, $\dfrac{\| m(kz)-k.m(z) \|_1}{|z|_f}, \dfrac{\| m(zk)-m(z) \|_1}{|z|_f} \to 0$ when $z\to \infty/\mathcal{G}$.
\end{enumerate}
\end{proposition}

\begin{proof}
(1) It suffices to prove that the set $\{ z\in H_{\Gamma}\rtimes G \mid |z|_f\leqq C\}$
is small relative to $\mathcal{G}$ for any positive integers $C$. 
But by the definition of $m(z)$ it easily follows that if $|(h_1h_2\cdots h_m, g)|_f\leqq C$ then the quantities $|\mathop{supp}h|$, $\min\{ |v|_{\Gamma}, |g^{-1}v|_{\Gamma}\}$, and $|h_i|_H$ are all smaller than $C$. 

Therefore we obtain
\[
\{ z\in H_{\Gamma}\rtimes G \mid |z|_f\leqq C\} \subset
A(\{ |\cdot |_H\leqq C\}, \{ |\cdot |_{\Gamma}\leqq C\}, C )
\]
and the conclusion follows by the assumption on $|\cdot|_{\Gamma}$ and $|\cdot|_H$. 

(2) It suffices to prove that, for any positive integer $C$, the set of all $z=(h,g)$ satisfying $|z|_f\leqq C |\mathop{supp} h|$ is small relative to $\mathcal{G}$. 
We claim that, if $|z|_f\leqq C |\mathop{supp} h|$, then
\[
|\mathop{supp} h| \leqq 4 |B(\Gamma,2C)|,
\] 
from which, combined with the assumption and (1), the conclusion follows. 
Here $B(\Gamma,R)=\{v\in V\Gamma \mid |v|_{\Gamma} \leqq R\}$ for a positive integer $R$. 

Suppose not. Then for every $g\in G$ one gets
\[
|B(\Gamma, 2C)\cup gB(\Gamma,2C)|< \dfrac{1}{2}|\mathop{supp}h|. 
\]
Using this inequality and the fact that $\min\{ |v|_{\Gamma}, |g^{-1}v|_{\Gamma}\}>2C$ for every $v\notin B(\Gamma, 2C)\cup gB(\Gamma,2C)$, the following inequality holds;
\begin{align}
|\mathop{supp}h| 
&\leqq |B(\Gamma, 2C)\cup gB(\Gamma,2C) | + |\mathop{supp}h \backslash (B(\Gamma, 2C)\cup gB(\Gamma,2C))| \\
&< \dfrac{1}{2}|\mathop{supp}h| + \dfrac{1}{2C}\sum_{v\in \mathop{supp}h \backslash B(\Gamma, 2C)\cup gB(\Gamma,2C)} \min\{ |v|_{\Gamma}, |g^{-1}v|_{\Gamma}\}\\
&\leqq \dfrac{1}{2}|\mathop{supp}h| + \dfrac{1}{2C}|(h,g)|_f \\
&\leqq |\mathop{supp}h|. 
\end{align}
This is a contradiction. 


(3) First, note that, by the proof of Lemma 2.12, the normal form of $xh$ should be one of the following form. 
\begin{itemize}
\item $xh_1h_2\cdots h_m$, when $\|xh\| = \| h\| +1$
\item $h_1\cdots h_{i-1}(xh_i)\cdots h_n$,  when $\|xh\| = \| h\|$
\item $h_1\cdots h_{i-1}h_{i+1}\cdots h_n$, when $\|xh\| = \| h\| -1$
\end{itemize}
In particular, $\mathop{supp}(xh)$ must be either $\{v\} \cup \mathop{supp}(h)$,  $\mathop{supp}(h)$, or $\mathop{supp}(h) \backslash \{ v \}$. 
Then, by Definition 4.4, 
\begin{align}
m(xz) = 
\begin{dcases}
\sum_{w\in \mathop{supp}(h)\cup \{ v \}} \min\{ |w|_{\Gamma}, |g^{-1}w|_{\Gamma}\} \delta_{w}+ \sum_{j=1}^m |h_i|_{H} \delta_{v_i} + |x|_H\delta_v,  
& (\|xh\| = \| h\| +1), \\
\sum_{w\in \mathop{supp}(h)} \min\{ |w|_{\Gamma}, |g^{-1}w|_{\Gamma}\} \delta_{w}+ \sum_{j=1, j\neq i}^m |h_i|_{H} \delta_{v_i} + |xh_i|_H\delta_v,  
& (\|xh\| = \| h\| ), \\
\sum_{w\in \mathop{supp}(h)\backslash \{ v \}} \min\{ |w|_{\Gamma}, |g^{-1}w|_{\Gamma}\} \delta_{w}+ \sum_{j=1, j\neq i}^m |h_i|_{H} \delta_{v_i}, 
& (\|xh\| = \| h\| -1).
\end{dcases}
\end{align}
In each case, by easy calculation we deduce that $\| m(xz)-m(z) \|_1 $ is smaller than $\min\{ |v|_{\Gamma}, |g^{-1}v|_{\Gamma}\} + |x|_H = |x|_f$. 
The right version follows similarly. 

(4) Note that $kz=(\sigma_k(h), kg)$. We see, by the definition of $m(z)$, that
\begin{align}
m(kz) &= \sum_{v\in k\mathop{supp}(h)} \min\{ |v|_{\Gamma}, |g^{-1}k^{-1}v|_{\Gamma}\} \delta_{v}+ \sum_{i=1}^m |h_i|_{H} \delta_{kv_i}, \text{and}
\\
k.m(z) &= \sum_{v\in \mathop{supp}(h)} \min\{ |v|_{\Gamma}, |g^{-1}v|_{\Gamma}\} \delta_{kv}+ \sum_{i=1}^m |h_i|_{H} \delta_{kv_i}. 
\end{align}
By taking the difference, it follows that 
\begin{align}
&\| m(kz)-k.m(z) \|_1 \\
\leqq 
&\| \sum_{v\in \mathop{supp}(h)} (\min\{ |kv|_{\Gamma}, |g^{-1}v|_{\Gamma}\} - \min\{ |v|_{\Gamma}, |g^{-1}v|_{\Gamma}\}) \delta_{kv} \|_1 \\
\leqq 
&\sum_{v\in \mathop{supp}(h)} | |kv|_{\Gamma} - |v|_{\Gamma} | \\
\leqq 
&|k|_G|\mathop{supp}h|
\end{align}
and the first inequality is proven. The second one can be proven similarly using 
$zk=(h,gk)$. 
\end{proof}

\end{document}
